\documentclass[a4paper,10pt]{article}
\usepackage[utf8]{inputenc}
\usepackage{graphicx}
\usepackage{subfigure}
\usepackage{fancyhdr}
\usepackage{listings}
\usepackage{ gensymb }


%opening

\pagestyle{fancy} %eigener Seitenstil
\fancyhf{} %alle Kopf- und Fußzeilenfelder bereinigen
\fancyhead[L]{SEP Abgabegespräch SS2015/16} %Kopfzeile links
\fancyhead[C]{} %zentrierte Kopfzeile
\fancyhead[R]{\includegraphics[width=1cm]{../../template/tugraz_bw}} %Kopfzeile rechts
\renewcommand{\headrulewidth}{0.4pt} %obere Trennlinie
\fancyfoot[C]{} %Seitennummer
\renewcommand{\footrulewidth}{0.4pt} %untere Trennlinie


\lstset{language=C++, basicstyle=\ttfamily\tiny}  

\begin{document}

\section*{Adressbuch}
Entwickle ein einfaches Programm, das Kontakte verwalten kann. Es gibt 2 Kontaktarten: Private Kontakte speichern neben dem Namen auch das Geburtsjahr der Person (als Ganzzahl). Geschäftliche Kontakte speichern neben dem Namen noch die E-Mail-Adresse. \\ Die Implementierung muss eine abstrakte Klasse beinhalten!

\begin{description}
 \item[add $<$type$>$ $<$name$>$ $<$parameter$>$] Erstellt einen Kontakt vom Typ $type$ ('private' oder 'business') und den angegebenen Parametern. 
 \item[list] Gibt alle erstellten Kontakte aus. 
\end{description}

Sollten die Paramter von $add$ nicht stimmen, soll der Befehl einfach ignoriert werden.
Bei der Eingabe von \texttt{quit} oder EOF soll das Programm beendet werden.


\subsection*{Beispiel}
\begin{verbatim}
> add private Peter 1994
> add private Simon asdf
> add business Dr.Hanser hanser@gmail.com
> list
Peter, 1994
Dr. Hanser, hanser@gmail.com
\end{verbatim}

\subsection*{Aufruf}
\texttt{./addressBook}

\subsection*{Bewertung}
\begin{itemize}
 \item Dokumentation und Programmierstil: 1 Punkt
 \item Einhalten der OOP Konzepte: 1 Punkt
 \item Struktur \& Korrektheit: 1 Punkt
\end{itemize}

\newpage


\end{document}
